% This is the Speculations to the Hodgkin-Huxley neuron model V2.0 book
% Last edited 2026.06.17


\chapter{Speculations\label{ch:Speculations}}

\section{Cell pressure\label{sec:CellPressure}}
The electrolyte contains dissociated molecules (free charges),
which, due to the repulsion between the ions, are located 
on the surfaces of the membrane and exert a pressure on the membrane surface.
We assume that the relationship between the amount of the charge on the membrane
and the pressure is linear, i.e., we conclude a pressure change from a charge change, and extrapolate it to the neuron in the resting state. 

\numericspanel{num:JohnstonConstants}
{Constants from Johnston\&Wu~\cite{JohnstonWuNeurophysiology:1995}, page~12}
{
\begin{tabular}{ll}

Neuron's surface 
$8*10^{-9}\ [m^2]$\\

\end{tabular}
}

\numericspanel{num:OtherConstants}
{Other measured constants}
{

Neuron's surface 
$8*10^{-9}\ [m^2]$\\

}


\physicspanel{phys:CellForceinField}{Estimating the force exerted on ions in electric field}
{ By using that in an electric field $\vec F = q \vec E$, one can calculate the force exerted on a charge, that enables one to calculate also the particle's impulse, pressure, and energy. For example, in the case of $Na^+$ rush-in, the electric force field is $10^7$ [V/m].
 Since($[N/C]==[V/m]$),
the value of the force exerted on a single ion, accelerating it across the ion channel, is
\begin{equation}
	F_{Na^+}= 10^7  [N/C] \times 1.602* 10^{-19} [C] = 1.6*10^{-12}\ [N] \label{eq:UnitForce}
\end{equation}

(Here we used data from Johnston\&Wu constants from \textbf{\hyperlink{num:JohnstonConstants}{Numerics panel~\ref{num:JohnstonConstants}}})
}

\physicspanel{phys:CellAccelerationinField}{Estimating the acceleration of ions in electric field}
{
The force in the ion channel was calculated in
\textbf{\hyperlink{phys:CellForceinField}{Physics Panel~\ref{phys:CellForceinField}}}
\noindent
and the mass of a $Na^+$ ion is $m_{Na} = 3.82\times 10^{-26}\ kg$.
It causes an acceleration
\begin{equation}
a=\frac{F_{Na^+}}{m_{Na^+}} = \frac{1.6*10^{-12}\ [N]}{ 3.82\times 10^{-26}\ [kg]} =
0.42 \times 10^{14}\ [\frac{m}{s^{2}}]
\label{eq:Physics-RushInAceleration}
\end{equation}
\noindent
The distance $d$ to travel (without collision) is the thickness of the neuronal membrane $5*10^{-9}\ nm$
and we assume a continuosly accelerating ion. Using $d=\frac{1}{2}\times a\times t^2$
\begin{equation}
t_{traverse} = \sqrt{\frac{2\times 5*10^{-9} [m]}{0.42\times 10^{14} [\frac{m}{s^{2}}]}} = 1.54\times 10^{-11}\ [s]
\end{equation}

The acceleration of an ion is unbelievably large. It is sufficiently large to keep the potential at the same value
over an extended object,
provided that the ions must follow a small and slow change, such as due to the
"leaking" current through the
\gls{AIS}.
Similarly, the ions can 'instantly' follow quick changes such as
a square wave gradient.
}


\physicspanel{phys:CellSpeedinField}{Estimating the speed of ions in electric field}
{

In \textbf{\hyperlink{phys:CellAccelerationinField}{Physics Panel~\ref{phys:CellAccelerationinField}}},
we estimated the acceleration and passage time of a $Na^+$ 
ion during a rush-in action.
In this period the ion is accelerated to the speed  $v=a*t$
\begin{equation}
v_{Na^+} = 0.42 \times 10^{14}\ [\frac{m}{s^{2}}]\ \times\ 1.54\times 10^{-11}  [s] = 600\ [\frac{m}{s}]
\end{equation}
(This speed is just a speed component along the channel. The thermal speed is about an order of magnitude higher, but not necessarily in the direction of the channel. The speed due to acceleration may convert in direction and magnitude due to a single collision, as soon as the diaphragm effect of the channel is over. The accelerated ion 'stops' as soon as
the accelerating voltage is over.)
See our other estimation in   
\textbf{\hyperlink{phys:RushinSpeed}{Physics Panel~\ref{phys:RushinSpeed}}}.
}

\physicspanel{phys:CellCurrentinChannel}{Estimating the current of ion channels}
{
Based on the time caculated in \textbf{\hyperlink{phys:CellAccelerationinField}{Physics Panel~\ref{phys:CellAccelerationinField}}},
we can attribute to an ion channel for this short  period,
provided that $10^3$ ions pass through a single ion channel, a current
\begin{equation}
I_{channel} = \frac{10^3 \times 1.60\times 10^{-19} [C]}{1.54\times10^{-11}[s]} = 10^{-5} A
\end{equation}
\noindent
It is in the range of several $\mu A$s, for a short period. It is distributed in the
spatiotemporal surface of a neuron, and contributes to a current of dozen of $pA$s.)

%The huge forces and accelerations mean that a potential change
%acts immediately. However, the force decays quickly. The ions start to move
%'instantly', but the charge carriers can move only with a limited speed,
%much below the interaction speed of
%\gls{EM}
%interactions. The effect can propagate only with that lower speed.
%
%See the case of axon: there exist a mechanical constraint that
%the ions cannot spread through the wall (they must keep the direction), and
%the ion package propagates as Equs.\ref{eq:Nernst-dVdt} and~\ref{eq:Nernst-dCdt} describe it.
}

\physicspanel{phys:CellPressure}{Estimating the pressure change due to rushin}
{
Using the force from \textbf{\hyperlink{phys:CellForceinField}{Physics panel~\ref{num:JohnstonConstants}}}, one can estimate how the pressure of the neural cell increases due to the rush-in changes at the beginning of the
\gls{AP}. 
As evidence shows, the local potential at the internal surface of the membrane is in the range of \SI{100}{\milli\volt} in the resting state and increases by $\Delta U=$ \SI{100}{\milli\volt}  in the transition state. This increase means a change in the force acting on an ion (see Eq.~(\ref{eq:UnitForce})) by  $1.6*10^{-12}\ [N]$.
When we assume $10^7$ rush-in ions and unchanged cell size, the total force acting on the membrane increases by $1.6*10^{-5}\ [N]$.
This change in force means on the neuron's $8*10^{-9}\ [m^2]$ surface  a pressure change 
\begin{equation}
	\Delta P = \frac{1.6*10^{-5}\ [N]}{8*10^{-9}\ [m^2]}
	= 2*10^{3}\quad \biggl[\frac{N}{m^2}\biggr]
	\label{eq:CellPressureChange}
\end{equation}
\noindent Again, the pressure can be calculated 
using electricity.

(Here we used data from Johnston\&Wu constants from \textbf{\hyperlink{num:JohnstonConstants}{Numerics panel~\ref{num:JohnstonConstants}}})
}

\physicspanel{phys:CellForceRushin}{Estimating the force measurable due to rushin}
{
The pressure change calculated in \textbf{\hyperlink{phys:CellPressure}{Physics panel~\ref{phys:CellPressure}}} enables to calculate
the force measurable dur to rushin.
	That pressure on a $10\times10~\mu m$ tooltip of the measuring device is converted to a force
	$2*10^{3}*100*10^{-12} = 200~pN$ force. The measured force value~\cite{MechanicalPropertiesNerves:2025} is about $600~pN$,
	so our estimation is in the correct range. 
	
	The measured pressure value is~\cite{PressureChangeActionPotentialMeasured:1980} $5\ \frac{dyn}{cm^2} = 0.5\ \frac{N}{m^2}$, which is an average value
	for some period. If we assume a $1~Hz$ frequency, and a $1~ms$
	period for the duration of the \gls{AP} peak, it means $1*10^{3}\, \bigl[\frac{N}{m^2}\bigr]$.
	So, the measured mechanical change values~\cite{NeuronalDeformation:2020} underpin our hypothesis that the increased pressure increases the neuron's size via its elasticity.
}

\physicspanel{phys:RushinSpeed}{Estimating the rush-in speed}
{
	A way to estimate the ions' rush-in speed is 
	that one assumes that the Stokes-Einstein speed describes
	the ions' speed (Fig.~3 in~\cite{HodgkinHuxley:1952}, quantitatively underpins the hypothesis, see~\cite{VeghStokesEinstein:2025}) both in axons and ion channels. If so, 
	the ratio of the potential gradients equals the ratio of the 
	speeds of ions. One can estimate that the $100~mV$ \gls{AP} on the 
	$50~\mu m$ \gls{AIS} generates a $2*10^4~\frac{V}{m}$ electrical gradient, and
	\cite{HodgkinHuxley:1952} measured $20~\frac{m}{s}$ for the speed of the \gls{AP}. By assuming that the electrical field across the membrane is $10^7~\frac{V}{m}$, one can estimate the speed at the exit of the ion channels as $10^4~\frac{m}{s}$.
	See our other estimation in
	\textbf{\hyperlink{phys:CellSpeedinField}{Physics Panel~\ref{phys:CellSpeedinField}}}
}


\physicspanel{phys:RushinMechanicalWork}{Estimating the rush-in mechanical work}
{
By using the result derived in \textbf{\hyperlink{phys:CellPressure}{Physics panel~\ref{phys:CellPressure}}},
	one can also estimate the work done when the rush-in charge
	extends the neuron, as the $\Delta P\times\Delta V$. We consider that $\Delta P$ remains constant, and we calculate the change of the volume as neuron's surface $10^{-8}\ m^2$ multiplied by 
	the change of the radius $\Delta r=10^{-9}\ [m]$ multiplied by the pressure $\Delta P$, that gives $E_{mech}=2*10^{-14}\ [J]$.
	If one assumes that the pressure is of entirely mechanical origin~\cite{NeuronalDeformation:2020}, one must assume that the 
	pressure is $10^6~\frac{N}{m^2}$, which is well above the measured value.
	}
	
\physicspanel{phys:RushinElectricalWork}{Estimating the rush-in electrical work}
{
	We can also estimate the electrical energy 
	from reset to threshold, see also~~\cite{EnergyNeuralCommunication:2021}
	\[
	E_{electr}=\frac{1}{2} C*(\Delta V)^2 = 	
	1.4*10^{-12}\ [J]
	\]
	 see also~\cite{EnergyNeuralCommunication:2021}.
}

\physicspanel{phys:RushinTime}{Estimating the rush-in time period}
{
	To estimate the time period required to change the pressure calculated in \textbf{\hyperlink{phys:CellPressure}{Physics panel~\ref{phys:CellPressure}}},
one can assume that the collision force of a $Na^+$ ion provides the pressure (the shock wave), and this can be calculated as $F = m\frac{\Delta v}{\Delta t}$. Given that Eq.~(\ref{eq:UnitForce})  provides the force, and one must assume that the ion entirely loses its 
	$10^4~to~10^5\frac{m}{s}$ speed, one concludes that it generates a shock wave in $\Delta t = 10^{-8}~to~10^{-7}~seconds$.
}

\physicspanel{phys:DissipatedEnergy}
{Estimating the dissipated energy of thermodynamical cycle}
{
We estimate the work done when the rush-in charge
	extends the neuron, as the $\Delta P\times\Delta V$. We consider that $\Delta P$ remains constant, and we calculate the change of the volume as neuron's surface $10^{-8}\ m^2$ multiplied by 
	the change of the radius $\Delta r=10^{-9}\ [m]$ multiplied by the pressure $\Delta P$,
that gives $\mathbf{E_{mech}=2*10^{-14}\ [J]}$, 
as we calculated in \textbf{\hyperlink{phys:RushinMechanicalWork}{Physics panel~\ref{phys:RushinMechanicalWork}}}.
%
\cite{MechanicalWaves:2015} reported that an \gls{AP} dissipates $Q_e$= \SI{0.25}{\micro\joule\per \meter\squared}, which on a \SI{8e-9}{\per\meter\squared}
surface means \SI{0.4e-14}{\joule}. 
As discussed in 
\textbf{\hyperlink{phys:RushinElectricalWork}{Physics panel~\ref{phys:RushinElectricalWork}}}, 
the electrical work is $\mathbf{E_{electr}=1.4*10^{-12}\ [J]}$.


The theoretically lowest possible energetic cost of information is $k_B T ln(2))$, numerically ca. \textbf{\SI{3e-21}{J}}
i.e., compared to the optimal limit set by physics, this efficiency value is $10^6$-$10^8$ times less energy efficient, in line with~\cite{EnergyNeuralCommunication:2021}. Nature protests:
the phenomenon is different from tha one biology claim. In the neuron there is no single state which can be set or reset.
}